\textbf{Tabel token cukup untuk sebagian besar topik berita (RQ1, RQ2).} StaticKD mempertahankan \val{statickd.test.pct_teacher} macro-F1 \eng{teacher} di set uji gabungan dan \val{statickd.dev.pct_teacher} di dev, dengan biaya inferensi dua sampai tiga orde besaran lebih rendah (\val{eff.statickd.speedup_onnx}$\times$ terhadap ONNX int8, \val{eff.statickd.speedup_pt}$\times$ terhadap PyTorch). Hal ini konsisten dengan temuan bahwa representasi tanpa urutan kata sudah kompetitif untuk klasifikasi topik~\cite{galke2022bow,iyyer2015dan}: nama lembaga, istilah olahraga, dan kosakata medis membawa sebagian besar sinyal topik. Sisa kesenjangan terkonsentrasi di dua tempat.
\begin{itemize}
  \item Label yang kosakatanya tumpang tindih dengan topik lain. Selisih terbesar terhadap \eng{teacher} ada pada \val{label.biggest_gap}. Label tersulit bagi StaticKD (\val{label.hardest}) juga termasuk yang tersulit bagi \eng{teacher} (Gbr.~\ref{fig:label}). \eng{Society} dan \eng{human interest} juga dikenal sulit bagi anotator manusia~\cite{kuzman2025llm}.
  \item Bahasa yang kurang terwakili di data distilasi. Pada bahasa dengan data distilasi cukup, StaticKD mencapai \val{grp.distil_tier.cukup.pct} skor \eng{teacher}; pada bahasa tanpa data distilasi, \val{grp.distil_tier.tidakada.pct}.
\end{itemize}
Dengan data distilasi yang sama, MiniLM-L6 menutup sebagian kesenjangan tersebut (selisih StaticKD $-$ MiniLM \val{sig.test.minilm_kd.delta}, $p$ \val{sig.test.minilm_kd.p}). DistilmBERT tidak menutupnya (selisih StaticKD $-$ DistilmBERT \val{sig.test.distilmbert_kd.delta}, $p$ \val{sig.test.distilmbert_kd.p}). Keduanya membutuhkan biaya satu sampai dua orde besaran lebih tinggi. Ini menunjukkan bahwa kualitas representasi multibahasa awal lebih menentukan daripada jumlah \eng{layer}: MiniLM berasal dari XLM-R-large, sedangkan DistilmBERT dari mBERT.

\textbf{Dari mana pengetahuan \eng{student} berasal (RQ3).} Ablasi memisahkan tiga sumber.
\begin{enumerate}
  \item \textbf{Data distilasi multibahasa yang dilabeli \eng{teacher}} adalah faktor terbesar. Tanpa data ini, macro-F1 uji turun sebesar \val{abl.emm_only.test_macro.delta}. Temuan ini sejalan dengan pandangan klasik bahwa nilai distilasi terletak pada \eng{transfer set} berlabel \eng{teacher}~\cite{bucilua2006model,ye2022sparse}.
  \item \textbf{Inisialisasi embedding yang selaras lintas bahasa.} Inisialisasi acak mengubah macro-F1 uji sebesar \val{abl.random_init.test_macro.delta}. Bukti yang lebih langsung datang dari varian \eng{compact}. Pada bahasa India yang hampir tidak ada di data distilasi (misalnya bn, pa, or), sebagian besar token tidak pernah diperbarui saat \eng{training}, sehingga embedding-nya tetap embedding awal potion. Meski begitu, StaticKD tetap mengklasifikasikan bahasa-bahasa ini dengan cukup baik (akurasi rata-rata \val{lang.vlowcov.static_acc}; Gbr.~\ref{fig:lang}). Ketika token-token tersebut dibuang oleh pemangkasan tokenizer, akurasinya jatuh ke \val{lang.vlowcov.compact_acc}. Jadi transfer ke bahasa tanpa data dibawa oleh baris embedding yang \emph{tidak pernah dilatih}, tetapi berada dalam ruang yang sama dengan baris yang dilatih. Kepala linear yang dipelajari dari bahasa kaya data berlaku juga untuk baris-baris ini.

  Ablasi embedding beku memperlihatkan sisi lain dari mekanisme yang sama. \eng{Fine-tuning} menggeser embedding token dari bahasa yang ada di data distilasi. Akibatnya kecocokan in-distribusi naik (dev \val{abl.frozen.dev_macro.delta} bila dibekukan), tetapi token bahasa lain tertinggal di ruang potion yang asli, sehingga kepala linear yang sudah menyesuaikan diri dengan ruang baru kurang cocok untuk token-token itu. Di bahasa tanpa atau dengan sedikit data distilasi, embedding beku lebih baik sebesar \val{ablgrp.frozen.tidakada.delta} dan \val{ablgrp.frozen.sedikit.delta}. Ini adalah bentuk pergeseran representasi (\eng{representation drift}) akibat \eng{fine-tuning} pada sebagian kosakata. Resep final tetap memakai \eng{fine-tuning} karena resep dipilih pada set \eng{development} yang berbahasa latih. Namun untuk penggunaan yang menargetkan bahasa tanpa data distilasi, varian beku atau regularisasi yang menahan embedding dekat ke potion adalah pilihan yang lebih aman. Keduanya tetap dapat dilipat menjadi tabel.
  \item \textbf{\eng{Soft label} memberi manfaat kecil atau tidak ada.} Label keras \eng{teacher} tidak berbeda signifikan dari \eng{soft label} dengan $\tau=1$ (dev \val{abl.hard.dev_macro.delta}, uji \val{abl.hard.test_macro.delta}). Sebaliknya, $\tau=2$ menurunkan kesepakatan \eng{held-out} dan macro-F1 uji, dan $\tau=4$ menurunkan semua metrik. Temuan ini sejalan dengan literatur \eng{capacity gap}~\cite{cho2019efficacy,zhang2025law,wang2026rckd}. \eng{Student} yang hanya berupa satu lapisan linear di atas rata-rata embedding tidak mampu merepresentasikan struktur antarkelas dalam ekor distribusi \eng{teacher} (\eng{dark knowledge})~\cite{zhao2022dkd,huang2022dist}. Distribusi yang lebih tajam ($\tau=1$) atau label keras memberi target yang dapat dicapai, sedangkan $\tau$ tinggi memaksa \eng{student} meniru skala dan bentuk logit yang di luar kapasitasnya~\cite{sun2024logit}.
\end{enumerate}
Kesimpulan praktisnya: untuk \eng{student} berkapasitas sangat rendah, sumber daya sebaiknya diarahkan ke cakupan \eng{transfer set} (bahasa, situs, topik) dan inisialisasi yang selaras, bukan ke penyetelan bentuk \eng{soft label}.

\textbf{Batas generalisasi lintas bahasa (RQ4).} Kinerja per bahasa terutama ditentukan oleh keterwakilan bahasa di data distilasi dan di pra-latih, bukan oleh sistem tulisan (Kruskal--Wallis $p$ \val{fac.kw.p}). Kelemahan terbesar ada pada bahasa Afrika yang berada di luar pra-latih XLM-R dan tidak ada di data distilasi (\val{lang.outxlmr.langs}). Pada bahasa-bahasa ini \eng{teacher} masih cukup kuat (akurasi \val{lang.outxlmr.teacher_acc}), mungkin berkat subword yang dipakai bersama dengan bahasa lain dan konteks kalimat, sedangkan StaticKD jauh lebih lemah (\val{lang.outxlmr.static_acc}). Solusi termurahnya tidak memerlukan perubahan arsitektur: tambahkan berita dalam bahasa-bahasa tersebut ke \eng{transfer set} dan biarkan \eng{teacher} melabelinya. Ablasi data dan korelasi kesepakatan dengan jumlah data ($\rho=\val{fac.data.rho}$) mendukung saran ini.

\textbf{\eng{Fidelity} tidak sama dengan akurasi.} Beberapa varian menggerakkan kesepakatan dengan \eng{teacher} dan akurasi terhadap label independen ke arah yang berlawanan.
\begin{itemize}
  \item Embedding beku menurunkan kesepakatan \eng{held-out} (\val{abl.frozen.heldout_acc.delta}) tetapi menaikkan macro-F1 uji (\val{abl.frozen.test_macro.delta}).
  \item Mengurangi data ke 130 ribu dokumen menurunkan kesepakatan (\val{abl.data130k.heldout_acc.delta}) tanpa perubahan signifikan di set uji (\val{abl.data130k.test_macro.delta}).
\end{itemize}
Sejalan dengan~\cite{stanton2021kd,ojha2023knowledge}, pemilihan \eng{student} sebaiknya memakai label yang independen dari \eng{teacher} dan set uji di luar distribusi latih, bukan \eng{fidelity} saja.

\textbf{Posisi terhadap penelitian sebelumnya.}
\begin{itemize}
  \item Dibanding \eng{Sparse Distillation}~\cite{ye2022sparse}, StaticKD mengganti tabel n-gram yang dilatih dari nol dengan embedding multibahasa pra-latih yang dilipat ke 17 logit per token. Model menjadi jauh lebih kecil dan dapat ditransfer ke bahasa yang tidak ada di \eng{transfer set}.
  \item Dibanding KD-LoRA dan PEFT~\cite{azimi2024kdlora,razuvayevskaya2024peft}, StaticKD menurunkan biaya inferensi, bukan biaya \eng{training}. Untuk \eng{pipeline} berita, biaya inferensi inilah yang dominan~\cite{luccioni2024power}.
  \item Kuantisasi int8 \eng{teacher} hanya memberi percepatan dengan faktor konstan~\cite{yao2022zeroquant}.
\end{itemize}

\textbf{Implikasi praktis.}
\begin{itemize}
  \item Untuk \eng{pipeline} \eng{real-time} di server tanpa GPU, StaticKD memproses satu juta artikel per hari dengan \val{eff.statickd.cores1M} \eng{core}, file \val{final.size}~MB, dan dependensi hanya \texttt{numpy} dan \texttt{tokenizers}.
  \item Probabilitas StaticKD lebih terkalibrasi daripada \eng{teacher} (ECE uji \val{ece.statickd.test} vs \val{ece.teacher.test}), sehingga dapat dipakai untuk mode \eng{cascade}. Dengan $\theta=0{,}8$, \val{casc.08.coverage} artikel diputuskan StaticKD dan sisanya dikirim ke \eng{teacher}. Akurasi keseluruhannya (\val{casc.08.accuracy}) setara dengan \eng{teacher} saja (\val{casc.101.accuracy}), dengan \eng{throughput} \val{casc.08.docs_per_sec} artikel per detik per \eng{core}, dibanding \val{eff.teacher_onnx_int8.docs_per_sec_1c} untuk \eng{teacher} saja. Mode ini cocok untuk jalur \eng{batch} yang membutuhkan akurasi setara \eng{teacher}.
  \item Varian \eng{compact} menurunkan RAM, tetapi hanya layak bila semua bahasa sasaran tercakup data distilasi.
\end{itemize}
